\section{A reference contract for continuity}\label{sec:C13:contract}
The following continuity contract is an original synthesis of the state and effect distinctions above. A resumed run should identify its session and task, load a validated event frontier, restore the authoritative goal and permissions, reconstruct the model-visible state through a versioned transformation, reconcile pending effects, and only then choose a new action. This ordering prevents the policy from treating uncertainty about external effects as if it were ordinary conversational ambiguity.

The contract does not require byte-identical future generation. A stochastic model can choose a different valid action after resumption. What should remain invariant is the set of obligations and the evidence about what has already happened. Reproducible sampling can help debugging, but operational continuity is not equivalent to deterministic decoding. A system that produces the same text while repeating an external action has reproduced the wrong aspect of the run.

Similarly, a session can migrate to a different model while retaining application continuity. That migration changes the policy interpreting the restored state and can alter how summaries, skills, or tool descriptions influence behavior. Version references make the change inspectable but do not prove compatibility. A migration gate should test representative obligations, stale-memory handling, and effect recovery under the new model. Chapter~\ref{ch:C17} analyzes this as a context-memory-policy compatibility problem.

For implementation, keep three stores logically distinct even if they share a database: immutable or policy-governed evidence, mutable task state, and derived model inputs. Evidence supports audit and reconstruction. Task state supports coordination and recovery. Derived inputs support efficient inference. Collapsing them into a single transcript makes every summarization operation responsible for persistence, authority, and learning provenance simultaneously. Separating the contracts makes failures attributable and allows different retention policies.

There is no implication that every agent application needs the strongest distributed guarantee. A read-only literature assistant can tolerate some duplicated searches, while a system dispatching externally consequential operations needs a more explicit effect model. The appropriate guarantee follows from the workload's failure costs and authority boundaries. The important scientific requirement is to state that guarantee precisely and test the intervals in which it could fail.

Durable state supplies continuity of evidence and obligations; it does not supply correct beliefs by itself. A faithfully restored session can still contain a mistaken inference or an obsolete semantic memory. The next chapter examines how memories are admitted, revised, and used, while Chapter~\ref{ch:C32} extends the operational analysis to deployment monitoring and rollback. Together these distinctions keep recovery claims bounded by the state that is actually preserved.
